% MNRAS submission version. Figures: copy the referenced PNGs from results/
% next to this file before compiling. Machine-readable tables:
% paper/table1_dispositions.csv, paper/table2_planet_limits.csv.
% Open items marked [PENDING]/[TODO]. AASTeX version retained as ms.tex.
\documentclass[fleqn,usenatbib]{mnras}

\usepackage[T1]{fontenc}
\usepackage{amsmath}
\usepackage{amssymb}
\usepackage{graphicx}

\newcommand{\msini}{m\sin i}
\newcommand{\mjup}{M_{\rm Jup}}
\newcommand{\mearth}{M_\oplus}

\title[Stability-derived true-mass limits for RV planets]{Stability-derived
upper limits on the true masses of non-transiting radial-velocity planets:
a homogeneous $N$-body survey}

\author[S. Ranjan]{Shaurya Ranjan$^{1}$\thanks{E-mail:
shauryaranjan911@gmail.com}\\
$^{1}$Independent Researcher}

\date{Draft \today}
\pubyear{2026}

\begin{document}
\label{firstpage}
\maketitle

\begin{abstract}
The radial-velocity method measures a planet's minimum mass, $\msini$, but
leaves the orbital inclination $i$, and therefore the true mass,
unconstrained. In a multi-planet system, long-term dynamical stability
supplies the missing bound. Because each planet's true mass scales as
$1/\sin i$, a system viewed near face-on contains heavier planets that
perturb one another more strongly, and above some inclination the
configuration no longer survives for the age of the system. The surviving
range of inclinations therefore sets an upper limit on each planet's true
mass. This argument has been applied to individual systems with a variety
of methods, but no uniform survey exists. We carry out a homogeneous
$N$-body stability survey of 56 multi-planet radial-velocity systems drawn
from the NASA Exoplanet Archive. For each system we sweep the
line-of-sight inclination over a fixed grid and sample all orbital
parameters within their published uncertainties. We validate the pipeline
on GJ~876, recovering the published stability constraint ($i > 25\degr$,
against the literature value $i \gtrsim 20\degr$). Under an isotropic prior
on orientation we derive 95 per cent credible upper limits on the true
mass of 123 planets in 53 systems. For 61 planets in 23 systems the limit
falls below $3.0\,\msini$, tighter than geometry alone provides (median
$2.68\,\msini$), and 25 giant planets are confined below the
deuterium-burning limit on dynamical grounds. We also document three ways
in which tabulated catalog parameters can be dynamically inconsistent, and
describe how the survey guards against each.
\end{abstract}

\begin{keywords}
celestial mechanics -- planets and satellites: dynamical evolution and
stability -- techniques: radial velocities -- methods: numerical --
catalogues
\end{keywords}

\section{Introduction} \label{sec:intro}

The radial-velocity method detects a planet through the reflex motion it
induces in its host star along the line of sight. This motion fixes the
product $\msini$, where $m$ is the planet mass and $i$ the orbital
inclination, but not $m$ itself. For a randomly oriented orbit the
correction from $\msini$ to $m$ is usually small: the median value of
$1/\sin i$ is 1.15, and 95 per cent of orientations give a true mass below
$3.2\,\msini$. For any individual planet, however, $\msini$ is a lower
bound rather than a measurement. This one-sided uncertainty propagates
into quantities built on radial-velocity masses. Mass functions and
occurrence rates are biased low; the boundary between planets and brown
dwarfs is ambiguous for any giant whose true mass could exceed $13\,\mjup$;
and inferences about the formation and composition of individual systems
often treat a lower limit as if it were the mass.

In a system of two or more planets, dynamics provides the upper bound that
the radial-velocity data lack. Consider a coplanar system tilted as a
rigid body towards a face-on orientation. Its internal dynamics are
unchanged, but every $\msini$ mass increases as $1/\sin i$. Heavier
planets exert stronger mutual perturbations, and for a sufficiently
face-on orientation the system becomes unstable on a time-scale far
shorter than its age. Such configurations cannot describe a system that is
observed to exist today, so the requirement of long-term survival places a
ceiling on each planet's true mass. This argument is well established for
individual systems. GJ~876 is the canonical example; stability has also
been used to constrain 47~UMa \citep{Gozdziewski2002}, HD~141399
\citep{Vogt2014}, and Barnard's star \citep{Basant2025}, and
\citet{Tamayo2021} developed it into a general method for
stability-constrained characterization, which they applied to Kepler-23.
What has been missing is a uniform treatment. \citet{LaskarPetit2017}
classified 131 systems by AMD stability but did not convert stability into
per-planet mass limits; \citet{Volpi2019} used first-order secular theory
to bound the mutual inclinations of ten radial-velocity systems; and
\citet{Qin2025} propose recovering inclinations from
von~Zeipel--Lidov--Kozai cycles in compact binaries. The remaining
literature consists of single-system analyses that adopt different
integrators, instability criteria, integration times, and statistical
definitions, so their mass limits cannot be compared and no
population-level statement can be drawn from them.

This paper presents a homogeneous survey that applies one pipeline, one
instability criterion, and one statistical definition to every qualifying
multi-planet radial-velocity system in the NASA Exoplanet Archive. Orbital
parameters are sampled within their published uncertainties, and the
pipeline is validated against systems whose stability constraints are
already known. We report 95 per cent credible upper limits on true mass
under an isotropic orientation prior, a set of systems flagged as
parameter-inconsistent, and three classes of catalog-parameter pitfall
that the dynamical treatment exposes (Section~\ref{sec:artifacts}).

\section{Sample selection} \label{sec:sample}

We draw candidate systems from the composite parameters table
(\texttt{pscomppars}) of the NASA Exoplanet Archive, accessed 2026
July 9--13. A system qualifies if it contains at least two planets with
measured periods and masses, has a measured stellar mass, and contains at
least one non-transiting planet whose tabulated mass is a minimum mass
($\msini$); these are the planets the method can constrain. Because mutual
perturbations weaken rapidly with orbital separation, we also require at
least one adjacent planet pair with a period ratio below 6. These criteria
yield 148 systems.

We rank each system by an analytic estimate of the mass factor at which it
should become unstable. For the most fragile adjacent pair, with masses
scaled by a common factor $f$ and separation measured in
eccentricity-reduced mutual Hill radii, we find the value $f_{\rm crit}$
at which the pair reaches the two-planet overlap boundary
$\Delta = 2\sqrt{3}$ \citep{Gladman1993}. Systems of more than two planets
destabilize at wider spacings, so $f_{\rm crit}$ overestimates the true
ceiling and is therefore a conservative selection criterion. We survey the
band $1.05 < f_{\rm crit} < 60$ with maximum eccentricity
$e_{\rm max} \le 0.75$ and inner-planet period $P_{\rm inner} \ge 1.5$~d,
which contains 56 systems. Systems with $f_{\rm crit} \approx 1$ are
excluded: this group is dominated by resonant and near-resonant pairs
whose survival depends on the relative orbital phases, information that the
composite catalog does not supply reliably
(Sections~\ref{sec:artifacts} and \ref{sec:discussion}).
Table~\ref{tab:dispositions} lists the disposition of all 148 systems: 53
surveyed, 3 flagged as parameter-inconsistent, and 92 excluded. The
excluded systems are resonance-protected (42), too weakly interacting for
stability to constrain the mass at any allowed value (48), or have an
inner period below 1.5~d (9), with some systems in more than one category.

\section{Methods} \label{sec:methods}

\subsection{Configuration sampling} \label{sec:sampling}

Tilting a coplanar system rigidly with respect to the line of sight
changes the inferred masses but not the internal dynamics: every $\msini$
mass scales as $1/\sin i$ while all mutual angles are preserved. We
therefore treat inclination as the single swept parameter. We sample $i$
on a 17-point grid from $90\degr$ (edge-on, masses at their minimum
values) to $5\degr$ (masses at 11.5 times the minimum). Planets with a
true mass known from a transit or from astrometry are held fixed. At each
grid point we integrate 20 independent realizations. In each realization
we draw the stellar mass and every planet's mass, period, eccentricity,
and argument of periastron from truncated normal distributions defined by
the symmetrized $1\sigma$ catalog uncertainties; where an uncertainty is
not tabulated we use the point value, and eccentricities are truncated to
$[0, 0.95]$. Angles the catalog does not report are drawn uniformly.
Orbital phases are randomized for systems whose parameters are assembled
from multiple references, and taken from the fitted time of periastron
only for systems with a single self-consistent reference, for the reasons
given in Section~\ref{sec:artifacts}.

\subsection{Integration and instability criterion}
\label{sec:integration}

We integrate each realization with the WHFast symplectic integrator
\citep{ReinTamayo2015} in the \textsc{rebound} package \citep{ReinLiu2012}.
The time-step is $dt = \min_p P_p (1-e_p)^{3/2} / 20$, one-twentieth of
the shortest perihelion passage time in the system, which keeps the
fastest pericentre well resolved. A realization is recorded as unstable if
any two bodies approach within the sum of their Hill radii (each body
carries its Hill radius as a collision radius) or if any body moves beyond
20 times the widest semi-major axis. Each realization is integrated for a
fixed number of inner-planet orbital periods rather than a fixed physical
time, so the computational cost per system does not depend on the absolute
scale of the system. We use two depths, $10^5$ and $10^6$ inner orbits;
Section~\ref{sec:depth} compares them.

We test the sensitivity of the limits to these choices on a representative
subset of six systems that span the survey's range of architectures.
Halving the time-step changes the 95 per cent credible limits by at most
0.03 in the mass factor. Reducing the collision radius to 0.5 and 0.25
times the Hill radius changes them by at most 0.08. Replacing exact
coplanarity with mutual inclinations drawn from Rayleigh distributions of
width $\sigma = 2\degr$ and $5\degr$ changes them by at most 0.4. None of
these shifts exceeds the grid spacing in inclination.

\subsection{Statistical definitions} \label{sec:stats}

For each system we construct a survival curve $S(i)$, the fraction of the
20 realizations at inclination $i$ that remain stable to $t_{\rm max}$. We
define the constraint relative to the edge-on value $S(90\degr)$, because
a system whose published parameters are already marginal can fail even at
edge-on orientations. A system with $S(90\degr) < 0.5$ is reported as a
parameter-consistency flag: its result describes the published parameters,
not a mass limit. For the remaining systems we report a 95 per cent
credible upper limit on the mass factor under an isotropic prior on
orientation (uniform in $\cos i$), taking the posterior density
proportional to $S(i)$; the quoted limit is the mass factor below which
95 per cent of the surviving posterior lies. With no dynamical information
this limit equals 3.20 for every planet, so values below 3.20 measure the
information the survey adds. For comparison with earlier single-system
work we also tabulate a threshold ceiling, defined as the smallest mass
factor at which $S$ falls below half its edge-on value.

\subsection{Catalog-parameter pitfalls guarded against}
\label{sec:artifacts}

During development we encountered three recurring ways in which tabulated
catalog parameters can be dynamically inconsistent. The survey guards
against each.

\begin{enumerate}
\item \textit{Eccentricity point estimates incompatible with survival.}
Compact systems are often fit with small, poorly constrained
eccentricities whose point values nonetheless destabilize the system when
integrated directly. For Barnard's star the tabulated values
$e = 0.03$--$0.08$ destroy the system within a few thousand years, whereas
eccentricities drawn from within the quoted uncertainties do not; the same
conclusion was reached independently by \citet{Basant2025}, whose fits
favour $e < 0.02$. We therefore sample eccentricities within their
uncertainties and never integrate the point values alone
(Section~\ref{sec:sampling}).
\item \textit{Composite parameters from incompatible fits.} The
\texttt{pscomppars} table assembles each parameter from the best available
source, so a single system can combine values from different fits. For
resonant systems this destroys the phase relationships on which stability
depends. We randomize the orbital phases of such systems and take fitted
phases only from a single self-consistent reference
(Section~\ref{sec:sampling}).
\item \textit{Fit epochs recorded as periastron times.} For some
references the tabulated time of periastron is the epoch of the fit, which
is identical for every planet and, if read literally, places all planets
at periastron simultaneously (GJ~876; \citealt{Rivera2010}). We reject a
system's periastron times when they are identical across its planets
before using any phase information.
\end{enumerate}

\section{Validation} \label{sec:validation}

We validate the pipeline on two systems chosen to bracket the two
dynamical regimes it encounters: GJ~876, whose survival depends on the
relative phases of a mean-motion resonance, and HD~45364, which is stable
independently of its phases. A third, external check appears in
Section~\ref{sec:catalog}, where the survey's mass ceilings for Barnard's
star are compared with the independent stability masses of
\citet{Byrne2026}.

\subsection{GJ 876: the phase-protected anchor} \label{sec:gj876}

GJ~876 is the standard test case for stability-constrained
characterization. The Archive's default solution is the coplanar fit of
\citet{Rivera2010} at $i = 59\degr$, which lists true masses and places
the 4:2:1 Laplace resonance at the centre of its libration. We reconstruct
the equivalent minimum-mass system by scaling the fitted masses by
$\sin 59\degr$ and fix the mean anomalies to the values in their Table~3.
For this reference the Archive's time-of-periastron column holds the epoch
of the fit rather than a per-planet periastron passage (pitfall 3 of
Section~\ref{sec:artifacts}), so the phases must be taken from the paper.
We then sweep a 13-point inclination grid with the phases held fixed,
using 5 realizations per grid point and 5000~yr ($\approx 10^6$ inner
orbits) per realization.

The test has three requirements (Fig.~\ref{fig:validation}). First, the
published configuration, which enters the sweep at mass factor
$1/\sin 59\degr = 1.17$, must survive; it does. Secondly, the recovered
ceiling must match the published constraint: survival is zero for all grid
inclinations at or below $25\degr$ and non-zero above, giving
$i > 25\degr$ ($m < 2.37\,\msini$), consistent with and slightly more
conservative than the literature value $i \gtrsim 20\degr$. Thirdly, the
same masses with randomized phases must not remain stable, since otherwise
resonant protection would play no role; survival falls to 12 per cent. The
pipeline thus reproduces both the resonant protection and the known
inclination limit.

\subsection{HD 45364: the phase-agnostic anchor} \label{sec:hd45364}

HD~45364 tests the opposite regime. Its two giant planets lie in a 3:2
commensurability with low eccentricities \citep{Correia2009}, and the
\citet{Li2022} solution is stable without help from resonant phases.
Initialized with the fitted phases, the published configuration survives
the full $10^5$~yr integration, and randomizing the phases leaves survival
unchanged at 100 per cent. The survival curve is correspondingly smooth
rather than sharpened by resonance: complete survival for $i \ge 45\degr$
falling to zero by $i \approx 8.5\degr$ (17-point grid, 10 realizations
per point). Under the isotropic prior this gives a 95 per cent credible
limit of $m < 2.29\,\msini$. The two anchors bracket the survey
population, one system constrained by its phases and one independent of
them, and the pipeline reproduces both without any per-system adjustment.

\subsection{Depth convergence} \label{sec:depth}

Because we integrate for a fixed number of orbits rather than to a
physical age, we must show that the reported limits reflect the dynamics
and not the integration length. We therefore run every system at two
depths, $10^5$ and $10^6$ inner orbits, using identical parameter draws
(the random seed is fixed by the system name), and compare the threshold
ceilings. A longer integration can only reveal additional instability, so
the ceilings should move in one direction. They do: among the 53 unflagged
systems, 9 ceilings tighten between the two depths, 5 are unchanged, none
loosen, and 7 systems that show no ceiling at $10^5$ orbits acquire one at
$10^6$ (Fig.~\ref{fig:depth}). This is the third successive generation of
the survey to show strictly one-directional behaviour, which is what
allows us to describe the reported limits as conservative: a longer
integration can only strengthen them (Section~\ref{sec:discussion}).

As a stronger test of the same property, we re-ran at $10^7$ inner orbits
the 16 systems whose ceilings had moved between $10^5$ and $10^6$ orbits.
Eight ceilings tighten further and eight are unchanged; none loosen, so
the one-directional behaviour persists at a tenfold-longer integration.

\section{Results} \label{sec:results}

\subsection{The catalogue} \label{sec:catalog}

Of the 56 systems in the survey band, 53 pass the edge-on consistency test
and yield 95 per cent credible mass limits for 123 non-transiting planets
(Table~\ref{tab:limits}; Fig.~\ref{fig:dumbbells}). For 61 planets in 23
systems the limit falls below $3.0\,\msini$, inside the value of 3.20 that
geometry alone guarantees, with a median of $2.68\,\msini$ among them.
Twenty-five giant planets whose face-on orientations would place them
above the deuterium-burning limit are instead confined below it at 95 per
cent credibility.

Several individual results illustrate the range. The four sub-Earth
planets of Barnard's star are limited to $2.18\,\msini$
(0.42--0.73\,$\mearth$), in agreement with the independent stability
masses of 0.19--0.84\,$\mearth$ reported by \citet{Byrne2026}. The three
giant planets of 47~UMa are limited to $2.35\,\msini$ and the four giant
planets of HD~141399 to $2.71\,\msini$. K2-18~c, the non-transiting
companion of the habitable-zone sub-Neptune K2-18~b, is limited to
$2.87\,\msini \approx 21.6\,\mearth$, which excludes a giant-planet mass.
For K2-18~c the isotropic prior is conservative: because K2-18~b transits,
only a small mutual inclination is needed for c to avoid transiting, so
near-edge-on orientations are favoured a priori \citep{Cloutier2017}.

The tightest limit in the catalogue shows that a small mass factor and a
planetary-mass classification are separate properties. The pair around
TYC~1422-614-1 is limited to $2.01\,\msini$, the strongest relative
constraint in the survey, but the minimum mass of the outer planet is
already $10\,\mjup$, so even this limit permits a mass of
$\approx 20\,\mjup$, above the deuterium-burning limit.

\subsection{Population statement} \label{sec:population}

Fig.~\ref{fig:cdf} shows the cumulative distribution of the 123 credible
limits. Half the sample (61 planets) is constrained below the geometric
value of 3.20, and the tightest fifth below $2.6\,\msini$. No limit
exceeds 3.20 by more than the discretization of the inclination grid; the
single nominal exception, TOI-1062~c at 3.22, is an unconstrained
posterior evaluated on the 17-point grid rather than a real excess. In
population terms, for closely spaced multi-planet radial-velocity systems
the high-mass tail implied by unconstrained orientation is largely
removed: the surviving configurations cluster near the edge-on masses, so
the true masses of these planets are typically close to their tabulated
minimum values.

\subsection{Parameter-consistency flags} \label{sec:flags}

Three systems fail the edge-on consistency test in more than half their
realizations: GJ~667~C, HIP~57274, and Kepler-20. These are statements
about the published parameters rather than mass limits: the tabulated
parameters, with their quoted uncertainties, do not admit a long-lived
configuration at any inclination. For GJ~667~C this is consistent with the
long-standing debate over the number of planets in the system; for the
other two it points to eccentricity solutions that are incompatible with
their own dynamics. In each case the appropriate next step is a joint
dynamical and radial-velocity refit.

\section{Discussion} \label{sec:discussion}

Three limitations bound the interpretation of the catalogue. First, the
inclination sweep is coplanar: each system is tilted as a rigid body, so
mutual inclinations enter only through the sensitivity tests of
Section~\ref{sec:integration}, which shift the limits by at most 0.4 in
the mass factor for dispersions up to $5\degr$. Strongly non-coplanar
systems, including those the von~Zeipel--Lidov--Kozai literature targets,
fall outside our assumptions, and their limits should be read as
conditional on approximate coplanarity. Secondly, $10^6$ inner orbits is
short compared with the Gyr ages of these systems. Every depth comparison
we have performed, across three generations of the survey, has moved the
limits monotonically towards tighter values or left them unchanged, so the
reported ceilings are conservative; we make no claim that a surviving
configuration is stable indefinitely. Thirdly, the inputs are catalog
parameters sampled within their symmetrized uncertainties, not fits to the
radial velocities themselves. The parameter-consistency flags identify
where this matters most; for the flagged systems the correct next step is
a joint refit rather than a mass limit.

The excluded resonant systems, with adjacent period ratios within a few
per cent of a small-integer commensurability, can be treated with the same
machinery validated on GJ~876, using single-reference parameters with
their fitted phases initialized at the observed libration. Extending the
survey to this group would add some of the most closely packed, and
therefore most strongly constrained, systems in the Archive.

The catalogue has several direct applications. Planets whose dynamical
ceiling lies close to their minimum mass have the narrowest permitted mass
range and are the best targets for astrometric confirmation with Gaia.
Planets whose ceiling falls below the deuterium-burning limit are
established as planetary-mass without additional observations, which
removes brown-dwarf contamination from the giant-planet mass function. The
parameter-consistency flags identify systems whose published solutions are
not supported by their own long-term dynamics and are candidates for
reanalysis.

\section*{Acknowledgements}

This work made use of the \textsc{rebound} $N$-body package
\citep{ReinLiu2012} with the WHFast integrator \citep{ReinTamayo2015}, and
of \textsc{numpy}, \textsc{pandas}, and \textsc{matplotlib}. It made use
of the NASA Exoplanet Archive, operated by the California Institute of
Technology under contract with the National Aeronautics and Space
Administration. Claude (Anthropic) was used to assist with code
development and manuscript preparation; all results were independently
verified by the author, who is responsible for their content.

\section*{Data availability}

The code and per-run records are available at
\url{https://github.com/Shaun2highh/dynamass}. The catalogues and per-run
records will be archived at Zenodo \textbf{[TODO: mint DOI at submission]}.
This work used the NASA Exoplanet Archive \texttt{pscomppars}
(DOI: 10.26133/NEA13) and \texttt{ps} (DOI: 10.26133/NEA12) tables,
accessed 2026 July 9--13.

% ---------------------------------------------------------------- figures

\begin{figure*}
\centering
\includegraphics[width=0.48\textwidth]{gj_876_validation.png}
\includegraphics[width=0.48\textwidth]{hd_45364_validation.png}
\caption{Validation anchors. Left: GJ~876 survival fraction versus
inclination with the \citet{Rivera2010} fit phases held fixed; survival is
zero at $i \le 25\degr$, recovering the literature constraint. Right:
HD~45364 \citep{Li2022}, a smooth, phase-agnostic survival curve (complete
survival at $i \ge 45\degr$, zero by $i \approx 8.5\degr$).}
\label{fig:validation}
\end{figure*}

\begin{figure*}
\centering
\includegraphics[width=0.9\textwidth]{floor_vs_ceiling_final.png}
\caption{Minimum masses ($\msini$) and 95 per cent credible dynamical
ceilings for the 123 constrained planets of the 53 unflagged systems,
sorted by credible mass factor. Flagged systems are excluded.}
\label{fig:dumbbells}
\end{figure*}

\begin{figure}
\centering
\includegraphics[width=\columnwidth]{depth_convergence_prod.png}
\caption{Depth convergence: per-system threshold ceilings at $10^5$ versus
$10^6$ inner-planet orbits with identical parameter draws. Ceilings
tighten (9 systems) or are unchanged (5); 7 systems acquire a ceiling only
at the greater depth; none loosen.}
\label{fig:depth}
\end{figure}

\begin{figure}
\centering
\includegraphics[width=\columnwidth]{m95_cdf_prod.png}
\caption{Cumulative distribution of the 123 95 per cent credible
mass-factor limits. The dashed line marks the isotropic-prior value (3.20)
that geometry alone guarantees; half the sample is constrained inside it.}
\label{fig:cdf}
\end{figure}

% ----------------------------------------------------------------- tables

\begin{table}
\centering
\caption{System dispositions (excerpt). Disposition of all 148 qualifying
systems: 53 surveyed, 3 flagged (edge-on survival below 0.5), and 92
excluded. The full machine-readable table is available as
\texttt{table1\_dispositions.csv}.}
\label{tab:dispositions}
\begin{tabular}{llrrr}
\hline
System & Disposition & $f_{\rm crit}$ & $S(90\degr)$ & $m_{95}$ \\
\hline
Barnard's star & surveyed & 1.60  & 0.55 & 2.18 \\
47 UMa         & surveyed & 1.34  & 0.55 & 2.35 \\
HD 141399      & surveyed & 3.05  & 0.85 & 2.71 \\
K2-18          & surveyed & 17.92 & 1.00 & 2.87 \\
GJ 667 C       & flagged  & 5.71  & 0.15 & -- \\
HIP 57274      & flagged  & 1.23  & 0.45 & -- \\
Kepler-20      & flagged  & 4.25  & 0.40 & -- \\
24 Sex         & excluded$^{a}$ & 1.00 & -- & -- \\
TOI-500        & excluded$^{b}$ & 32.86 & -- & -- \\
\hline
\end{tabular}
\\[2pt]
\raggedright
$^{a}$Resonance-protected / at the stability limit
($f_{\rm crit} \approx 1$).
$^{b}$Inner period below 1.5~d.
\end{table}

\begin{table}
\centering
\caption{Per-planet 95 per cent credible true-mass limits (ten tightest of
123, sorted by mass factor). The full machine-readable table
(\texttt{table2\_planet\_limits.csv}) additionally lists masses in
$\mjup$, the threshold ceiling, and the inclination floor.}
\label{tab:limits}
\begin{tabular}{llrrr}
\hline
System & Planet & $\msini$ ($\mearth$) & $m_{95}$ & $m_{95}$ ($\mearth$) \\
\hline
TYC 1422-614-1 & b & 794.5 & 2.01 & 1600 \\
TYC 1422-614-1 & c & 3178  & 2.01 & 6400 \\
Barnard's star & b & 0.299 & 2.18 & 0.65 \\
Barnard's star & c & 0.335 & 2.18 & 0.73 \\
Barnard's star & d & 0.263 & 2.18 & 0.57 \\
Barnard's star & e & 0.193 & 2.18 & 0.42 \\
HD 95089       & b & 400.4 & 2.31 & 923  \\
HD 95089       & c & 1096  & 2.31 & 2527 \\
HIP 111909     & b & 384.5 & 2.33 & 897  \\
HIP 111909     & c & 257.4 & 2.33 & 601  \\
\hline
\end{tabular}
\end{table}

% ------------------------------------------------------------- references

\begin{thebibliography}{}
% All entries verified against Crossref DOI metadata, 2026-07-22.
\bibitem[Basant et al.(2025)]{Basant2025} Basant R., Luque R., Bean
J.~L., et al., 2025, ApJL, 982, L1
\bibitem[Byrne et al.(2026)]{Byrne2026} Byrne X., Guimond C.~M., Bonsor
A., et al., 2026, MNRAS, accepted (arXiv:2606.23310)
\bibitem[Cloutier et al.(2017)]{Cloutier2017} Cloutier R.,
Astudillo-Defru N., Doyon R., et al., 2017, A\&A, 608, A35
\bibitem[Correia et al.(2009)]{Correia2009} Correia A.~C.~M., Udry S.,
Mayor M., et al., 2009, A\&A, 496, 521
\bibitem[Gladman(1993)]{Gladman1993} Gladman B., 1993, Icarus, 106, 247
\bibitem[Go\'zdziewski(2002)]{Gozdziewski2002} Go\'zdziewski K., 2002,
A\&A, 393, 997
\bibitem[Laskar \& Petit(2017)]{LaskarPetit2017} Laskar J., Petit A.~C.,
2017, A\&A, 605, A72
\bibitem[Li et al.(2022)]{Li2022} Li Z., et al., 2022, AJ, 164, 163
\bibitem[Qin et al.(2025)]{Qin2025} Qin Z., Liu S.-F., Ma B., 2025, ApJL,
988, L37
\bibitem[Rein \& Liu(2012)]{ReinLiu2012} Rein H., Liu S.-F., 2012, A\&A,
537, A128
\bibitem[Rein \& Tamayo(2015)]{ReinTamayo2015} Rein H., Tamayo D., 2015,
MNRAS, 452, 376
\bibitem[Rivera et al.(2010)]{Rivera2010} Rivera E.~J., Laughlin G.,
Butler R.~P., et al., 2010, ApJ, 719, 890
\bibitem[Tamayo et al.(2021)]{Tamayo2021} Tamayo D., Gilbertson C.,
Foreman-Mackey D., 2021, MNRAS, 501, 4798
\bibitem[Vogt et al.(2014)]{Vogt2014} Vogt S.~S., et al., 2014, ApJ, 787,
97
\bibitem[Volpi et al.(2019)]{Volpi2019} Volpi M., Roisin A., Libert
A.-S., 2019, A\&A, 626, A74
\end{thebibliography}

\bsp
\label{lastpage}
\end{document}
